\documentclass[a4paper]{article}
% Español (México) con XeLaTeX: fontspec para fuentes Unicode, polyglossia
% para la división silábica y los nombres del documento en la variante mexicana.
\usepackage{fontspec}
\usepackage{polyglossia}
\setdefaultlanguage[variant=mexican]{spanish}
% Los nombres chinos van como «pinyin (original)»: xeCJK da a los caracteres
% Han su propia fuente; sin ella XeLaTeX los omite con un aviso.
% FandolSong no cubre algunos caracteres de nombres (U+73FA); WenQuanYi Zen Hei sí.
\usepackage[AutoFallBack=true]{xeCJK}
\setCJKmainfont{FandolSong-Regular.otf}
\setCJKfallbackfamilyfont{\CJKrmdefault}{WenQuanYi Zen Hei}
% polyglossia no traduce los nombres de listings.
\AtBeginDocument{\providecommand{\lstlistingname}{}\renewcommand{\lstlistingname}{Listado}%
  \providecommand{\lstlistlistingname}{}\renewcommand{\lstlistlistingname}{Índice de listados}}
\usepackage{amsmath, amssymb}
\usepackage{graphicx}
\usepackage[margin=2.35cm]{geometry}
\usepackage[most]{tcolorbox}
\usepackage{booktabs}
\usepackage{tabularx}
\usepackage{longtable}
\usepackage{array}
\usepackage{float}
\usepackage{xcolor}
\usepackage{hyperref}
\usepackage{fancyhdr}
\hypersetup{colorlinks=true, linkcolor=blue!55!black, urlcolor=blue!60!black}
\graphicspath{{./}{figures/}}
\setlength{\parskip}{0.55em}
\setlength{\parindent}{2em}
\setlength{\emergencystretch}{3em}
\renewcommand{\arraystretch}{1.18}
\newtcolorbox{knowledgebox}[1]{enhanced, colback=blue!5!white, colframe=blue!70!black, colbacktitle=blue!70!black, coltitle=white, fonttitle=\bfseries, title={#1}, sharp corners, breakable}
\newtcolorbox{importantbox}[1]{enhanced, colback=yellow!10!white, colframe=yellow!75!black, colbacktitle=yellow!75!black, coltitle=black, fonttitle=\bfseries, title={#1}, sharp corners, breakable}
\newtcolorbox{warningbox}[1]{enhanced, colback=red!5!white, colframe=red!70!black, colbacktitle=red!70!black, coltitle=white, fonttitle=\bfseries, title={#1}, sharp corners, breakable}
\pagestyle{fancy}
\fancyhf{}
\fancyfoot[C]{\thepage}
\fancyfoot[R]{\footnotesize\color{gray}\href{https://github.com/hqhq1025/ai-course-notes}{AI Course Notes}}
\renewcommand{\headrulewidth}{0pt}
\renewcommand{\footrulewidth}{0pt}
\newcommand{\videourl}{https://www.youtube.com/watch?v=ruVJ_5dObxs}
\newcommand{\repourl}{https://github.com/hqhq1025/ai-course-notes}

\begin{document}
\begin{titlepage}
\centering
{\Large Notas didácticas de entrevista a fondo\par}
\vspace{1.0cm}
{\huge\bfseries Productos AI Native, diseño de contexto y apuestas por los Agent\par}
\vspace{0.5cm}
{\Large Zhang Xiaojun conversa con Zhang Yueguang (张月光)\par}
\vspace{0.35cm}
{\large Elaborado a partir del video público de Zhang Xiaojun Podcast, la transcripción hecha con GPU, la estructura de capítulos y diagramas conceptuales propios\par}
\vspace{0.3cm}
{\large 2026-01-22\par}
\vspace{0.85cm}
\includegraphics[width=0.88\textwidth,height=0.42\textheight,keepaspectratio]{cover.jpg}\par
\vfill
\begin{tcolorbox}[width=0.92\textwidth, colback=black!2!white, colframe=black!55, sharp corners]
\textbf{Título del video}: 130. Primera entrevista de Zhang Yueguang tras dos años de emprendimiento: Miaoya (妙鸭) no es un producto AI Native, del flujo de trabajo al diseño de contexto, One Way Door y los juegos otome\par
\textbf{Canal del video}: Zhang Xiaojun Podcast\par
\textbf{Duración del video}: 03:14:15\par
\textbf{Enlace del video}: \href{\videourl}{\nolinkurl{\videourl}}\par
\textbf{Nota sobre los materiales}: YouTube no ofrece subtítulos disponibles; el texto se elaboró a partir de la transcripción con faster-whisper large-v3 en CUDA, los capítulos del video, la portada y diagramas didácticos propios. Las tomas fijas de la entrevista no se reutilizan en el cuerpo del texto.\par
\textbf{Página del proyecto}: \href{\repourl}{\nolinkurl{\repourl}}\par
\end{tcolorbox}
\end{titlepage}

\tableofcontents
\newpage
\section{Introducción: de un producto fenómeno a una nueva comprensión de AI Native}

El punto de partida de esta entrevista es Miaoya (妙鸭). En 2023, Zhang Yueguang (张月光) creó en Alibaba (阿里) Miaoya, un producto de retratos fotográficos con IA que se volvió un fenómeno; sin embargo, dos años después afirma que Miaoya no es un producto AI Native. Esta afirmación es el giro conceptual más importante de todo el episodio. No niega el valor de Miaoya, sino que distingue entre «usar tecnología de IA» y «que la IA reescriba el paradigma del producto».

Esta sección establece primero el hilo conductor de todo el texto: Zhang Yueguang, un gerente de producto formado en la era de internet, pasó por grandes empresas tecnológicas, un primer emprendimiento, Miaoya, su salida de Alibaba, la fundación de Muyan Zhiyu (沐言智语) y la exploración de Xingmian (星眠) y Dokie, y finalmente reduce el problema a dos juicios. Primero, lo esencial de un producto AI Native no es si usa IA, sino el paso del diseño de flujos al diseño del contexto. Segundo, la apuesta más valiosa para la siguiente etapa es el Agent, en especial un amigo de IA capaz de superar los límites de la capacidad individual.

\begin{importantbox}{Tesis central del episodio}
AI Native no se define simplemente como «lo que no puede hacerse sin IA». Una definición de nivel superior es esta: tanto la entrada del usuario como la salida del modelo se vuelven abiertas, de modo que el gerente de producto ya no puede enumerar todos los flujos, sino que debe diseñar el contexto, las herramientas, la combinación de modelos, la retroalimentación y la percepción del producto en la mente del usuario.
\end{importantbox}

\begin{knowledgebox}{Nota sobre la estrategia visual}
Este video es una entrevista con encuadre fijo, sin slides, pizarrón ni demostraciones de producto. Conforme al estándar de notas de pódcast, el cuerpo del texto no repite fotogramas de las personas; la portada sirve para identificar la fuente, y el cuerpo del texto usa diagramas conceptuales, como el paso del flujo al contexto, One Way Door, la organización de productos de IA y el ciclo de retroalimentación del Agent, para explicar el contenido.
\end{knowledgebox}

\subsection{Resumen de la sección}

Este episodio no es solo una «revisión retrospectiva de Miaoya». Lo que realmente discute es cómo un gerente de producto de internet atraviesa el cambio de paradigma de la IA: de los flujos, los botones y los embudos, hacia el contexto, las capacidades del modelo, las acciones del Agent y la percepción del usuario.
\section{La formación de un gerente de producto: el juego de las grandes empresas y la «pendiente larga con nieve espesa»}

La introducción anterior planteó la gran pregunta de AI Native; esta sección vuelve primero a la persona. Lo que Zhang Yueguang (张月光) piensa de los productos de IA no surge de la investigación en modelos, sino de más de diez años como gerente de producto, dos emprendimientos y el juego organizacional de las grandes empresas.

La trayectoria de Zhang Yueguang es muy variada: SaaS, Alipay (支付宝), ByteDance (字节), su primer emprendimiento, Youku de Alibaba (阿里优酷), Miaoya (妙鸭) y su segundo emprendimiento. En la entrevista define los años de los 35 a los 45 como una etapa nueva, en la que quiere hacer una misma cosa con un mismo grupo de personas. Eso explica por qué dejó Alibaba: no por el éxito o el fracaso de un producto en particular, sino porque quería librarse de los cambios incontrolables de una gran empresa y volver a organizar a las personas y el trabajo.

\subsection{Las dos batallas del juego en las grandes empresas}

Divide el avance dentro de una gran empresa, como quien sube de nivel venciendo enemigos, en dos tipos: batallas externas y batallas internas. Las batallas externas consisten en ganar el mercado, ganar usuarios y ganarles a los competidores; las internas, en ganarse a quienes te rodean, ganar en el sistema de evaluación y ganar en la asignación de recursos. En la fase de crecimiento de una gran empresa, el elevador va hacia arriba y abundan las batallas externas; cuando la industria entra en un periodo de vacío, las batallas internas se vuelven más atractivas y es fácil que los mecanismos de la organización arrastren a las personas de vuelta a la competencia interna.

\begin{warningbox}{El costo oculto de las batallas internas}
Las batallas internas pueden traer ascensos, pero poco a poco desplazan la atención del gerente de producto: de «qué les pasa a los usuarios y al mundo exterior» a «cómo me ve tal o cual persona dentro de la organización». La nueva oportunidad que ofrece la IA es justamente reabrir el juego externo.
\end{warningbox}

\subsection{El hilo emocional del primer emprendimiento}

El producto Yuanyin (元音), de su primer emprendimiento, le hizo sentir por primera vez de forma directa a «usuarios de carne y hueso». Los usuarios conversaban en los grupos, pagaban y, cuando se cerró el servicio, escribieron largos textos de despedida; algo muy distinto de la cifra abstracta de 5 millones de usuarios nuevos que se ve en el panel interno de una gran empresa. Esta experiencia continuó después en Xingmian (星眠): el juego otome con IA no surgió de la nada, sino de su interés de largo plazo, desde temprano, por Live2D, la compañía, los personajes y las usuarias.

\begin{knowledgebox}{De las cifras de usuarios a la sensación de usuarios vivos}
Los datos de una gran empresa entrenan el sentido de escala, pero no necesariamente el juicio emocional sobre un producto. La retroalimentación de los usuarios de un producto pequeño permite al gerente de producto ver a personas concretas, gustos concretos y decepciones concretas, y eso influye en su juicio posterior sobre la compañía mediante IA y los amigos de IA.
\end{knowledgebox}

\subsection{Resumen de la sección}

El impulso de fondo de Zhang Yueguang es volver del juego interno al juego externo, de los indicadores abstractos a los usuarios concretos, y pasar de los proyectos de corto plazo a compañeros y a una empresa de largo plazo. Este antecedente determina su preferencia posterior por One Way Door, los amigos de IA y los Agent.
\section{Miaoya: un producto de internet que usa bien la IA}

Esta sección pasa de la trayectoria personal a un caso de producto. Entender Miaoya (妙鸭) es la primera puerta para entender toda la entrevista: es el caso de éxito más conocido de Zhang Yueguang (张月光) y también el contraejemplo que más tarde usó para reflexionar sobre «qué es realmente AI Native».

El éxito de Miaoya proviene de varias decisiones de producto: la tecnología de imagen ya era relativamente madura en ese momento; no hacer generación de imágenes de propósito general al estilo de Midjourney; elegir un nicho suficientemente vertical, el retrato fotográfico; evitar la larga cadena del comercio electrónico y optar por un negocio de retratos dirigido al consumidor (2C) al que se le puede poner precio; y reducir el objetivo de resultado a «realista, parecido, bello». Todas estas decisiones son excelentes, pero Zhang Yueguang concluyó después que Miaoya seguía respondiendo a una mentalidad de producto de internet.

\begin{figure}[H]
\centering
\includegraphics[width=0.94\textwidth]{miaoya-not-ai-native.png}
\caption{Por qué Miaoya no es AI Native: entrada fija, cadena de procesamiento fija, salida fija.}
\end{figure}

\begin{knowledgebox}{Lectura de la figura: la estabilidad de Miaoya proviene de limitar los grados de libertad}
Los tres segmentos de la figura son entrada fija, cadena de procesamiento fija y salida fija. El usuario sube unas 20 fotos, el sistema entrena un modelo personal y luego produce los retratos mediante plantillas fijas y una cadena de generación fija. Usa muy bien la capacidad de la IA, pero sigue obteniendo resultados estables a cambio de limitar la libertad del usuario.
\end{knowledgebox}

\subsection{Realista, parecido, bello}

Antes se dijo que la cadena de procesamiento de Miaoya es fija, pero una cadena fija no equivale a una cadena sencilla. Su verdadera dificultad está en lograr, para un escenario muy acotado, una función de efecto lo bastante sólida.

El estándar de resultado de Miaoya consta de tres elementos. Realista: que no se note que lo generó una IA; parecido: que se parezca lo suficiente al propio usuario, de modo que lo que se comparte tenga que ver con «yo»; bello: que se vea un poco mejor que uno mismo en la realidad. Este triángulo explica por qué se volvió un éxito masivo: no exhibe la tecnología de IA, sino que la convierte en una experiencia de consumo intensa.

\begin{importantbox}{La lección de producto de Miaoya}
Un éxito masivo de IA no necesariamente proviene del modelo más general; puede provenir de una función de efecto bien definida dentro de un escenario vertical. La función de efecto de Miaoya es justamente «realista, parecido, bello».
\end{importantbox}

\subsection{Por qué sigue siendo un producto de internet}

La función de efecto explica por qué Miaoya tuvo éxito; esta sección explica por qué, aun así, no es AI Native. Lo decisivo no es si el resultado deslumbra, sino si el producto sigue dependiendo de un flujo predefinido.

En esencia, Miaoya sigue siendo un diseño de flujo: el usuario sube fotos siguiendo pasos predefinidos, elige una plantilla y el sistema genera el resultado a lo largo de una cadena de proceso fija. La entrada del usuario no es abierta y la salida tampoco. El gerente de producto puede diseñar el flujo completo, los ingenieros lo implementan tal como está definido y la experiencia es controlable. Es un triunfo de la metodología de producto de internet, no el despliegue completo de la metodología AI Native.

\subsection{Resumen de la sección}

Miaoya muestra que usar bien la capacidad de la IA no equivale a que la IA haya reescrito el paradigma de producto. Es un excelente «producto de internet con IA añadida», pero no es AI Native en el sentido en que Zhang Yueguang lo entendió después.
\section{Del diseño de flujos al diseño de contexto}

La metodología de producto más importante de este episodio es el paso del flow design al context design. Durante los últimos 20 años, el paradigma básico de los productos de internet ha sido este: el gerente de producto define de antemano la ruta del usuario, las entradas y las salidas; el diseñador dibuja la interfaz, y el ingeniero implementa el flujo. En los productos AI Native, la entrada del usuario es abierta y la salida del modelo también lo es. Sobre todo en la era de los Agent, el modelo además llama herramientas, genera archivos y ejecuta acciones, por lo que el flujo no puede enumerarse de forma exhaustiva.

\begin{figure}[H]
\centering
\includegraphics[width=0.96\textwidth]{flow-to-context-design.png}
\caption{Del diseño de flujos al diseño de contexto: de grados de libertad limitados a una entrada abierta y una salida abierta.}
\end{figure}

\begin{knowledgebox}{Lectura de la figura: en AI Native cambió el objeto del diseño}
A la izquierda, el diseño de flujos se ocupa de botones, páginas, embudos y salidas fijas; a la derecha, el diseño de contexto se ocupa del prompt, el context, las herramientas, la combinación de modelos, las restricciones y la retroalimentación. El flujo del producto sigue siendo importante, pero debe deducirse a partir de las capacidades del modelo y del efecto del contexto, en lugar de dibujar primero el flujo completo y esperar que el modelo se ajuste a él.
\end{knowledgebox}

\subsection{Por qué falla la colaboración tradicional}

La figura anterior muestra la diferencia de paradigma; aquí se aplica a la colaboración del equipo. AI Native no es un cambio de mentalidad exclusivo del gerente de producto, sino un cambio en el orden en que colaboran los equipos de producto, diseño, ingeniería y modelos.

La colaboración tradicional de un equipo es lineal: el gerente de producto escribe los requisitos, el diseñador hace los diseños y el ingeniero los construye. En un producto AI Native, la parte intermedia más importante puede ser el context, que el usuario no ve: qué materiales se le dan al modelo, cómo se organiza la tarea, qué herramientas se llaman, cómo se evalúa la salida y cómo se incorpora la retroalimentación del usuario a la siguiente ronda. Si estos elementos no se hacen funcionar primero de principio a fin, cuanto más refinado sea el flujo de la UI, más equivocado puede estar.

\begin{warningbox}{El riesgo de empezar por el flujo}
Si primero se diseñan por completo las páginas, los botones y las rutas del usuario, y después se conecta el modelo, es muy probable descubrir que la salida real del modelo no coincide con el flujo previsto. Un producto AI Native debe validar primero el ciclo cerrado de capacidades y después diseñar el flujo externo.
\end{warningbox}

\subsection{Resumen de la sección}

El paso del diseño de flujos al diseño de contexto es la brecha conceptual decisiva para el gerente de producto de internet que se mueve hacia los productos de IA. El gerente de producto ya no es solo el diseñador de la ruta del usuario, sino que debe convertirse en quien organiza las capacidades del modelo, el contexto, la retroalimentación y el modelo mental del usuario.
\section{One Way Door: señales de un producto acertado a largo plazo}

Tras la discusión sobre Miaoya (妙鸭) y el diseño del contexto, hace falta un criterio para seleccionar productos. Zhang Yueguang (张月光) usa One Way Door para juzgar si vale la pena apostar a largo plazo por algo.

A Zhang Yueguang le gustan los productos One Way Door: una vez que el usuario adopta la nueva solución, sabe que ya no puede volver a la anterior. Cursor y Claude Code, frente a los IDE tradicionales, son un ejemplo típico. Un One Way Door no necesariamente ofrece una experiencia perfecta desde el principio, pero deja ver que en el futuro formará una ventaja integral.

\begin{figure}[H]
\centering
\includegraphics[width=0.92\textwidth]{one-way-door-product.png}
\caption{Criterio de producto One Way Door: solución anterior, nueva solución, nueva normalidad.}
\end{figure}

\begin{knowledgebox}{Lectura de la figura: la puerta de un solo sentido no es un wow moment}
En la figura, la puerta de un solo sentido se ubica entre la solución anterior y la nueva normalidad. Un wow moment solo deslumbra, y el usuario puede volver pronto al flujo anterior; la puerta de un solo sentido, en cambio, modifica el comportamiento predeterminado del usuario. Su participación de mercado puede empezar siendo pequeña, pero crece de forma sostenida con el tiempo.
\end{knowledgebox}

\subsection{Por qué un Agent puede ser una puerta de un solo sentido}

Si un producto Agent realmente permite al usuario superar los límites de su capacidad, no solo aumenta la eficiencia, sino que cambia la idea que el usuario tiene de «qué puedo hacer». Zhang Yueguang explica que Dokie entra por ahora a través de las PPT, pero que su verdadero objetivo no son las PPT, sino un amigo de IA que le ayude a superar los límites de su capacidad personal.

\begin{importantbox}{Prueba de producto de la puerta de un solo sentido}
Para saber si vale la pena desarrollar un producto de IA a largo plazo, pueden plantearse estas preguntas: después de usarlo, ¿no quiero volver al método anterior? ¿Resuelve una necesidad existente, pero de una manera nueva que ofrece una ventaja integral? ¿Se volverá cada vez más potente conforme avancen los modelos?
\end{importantbox}

\subsection{Resumen de la sección}

One Way Door es el filtro de productos de este episodio. Miaoya fue un gran éxito viral, pero no necesariamente una puerta de un solo sentido; un Agent, si logra superar de manera estable los límites de capacidad, puede convertirse en una puerta acertada a largo plazo.
\section{Lista de evaluación de producto: del Wow a la Retention}

Las secciones anteriores trataron Miaoya (妙鸭), AI Native y One Way Door, pero el juicio sobre un producto necesita además una lista más aplicable. En la entrevista, Zhang Yueguang (张月光) separa varias veces lo «vistoso» y lo «deslumbrante» de lo «correcto a largo plazo»: un producto puede ser muy wow y no ser una puerta de un solo sentido; puede tener un auge a corto plazo y no retener a los usuarios a largo plazo; puede ser técnicamente muy sólido y no ocupar un lugar en la mente del usuario.

\subsection{Evaluación en cuatro niveles}

\begin{longtable}{p{0.20\textwidth}p{0.36\textwidth}p{0.35\textwidth}}
\toprule
Nivel & Pregunta que hay que hacer & Lo que enseñan Miaoya / Dokie \\
\midrule
Viabilidad técnica & ¿Puede esta generación de modelos producir resultados de forma estable? & Miaoya eligió la tecnología de imagen relativamente madura en ese momento; Dokie debe comprobar si el Agent puede completar de verdad la creación y la expresión. \\
Función de efecto & ¿Por qué el usuario siente que es bueno? & En Miaoya es «real, parecido, bello»; en Dokie es «superar los límites de la propia capacidad». \\
Puerta de un solo sentido & Después de usarlo, ¿el usuario ya no puede volver al método anterior? & Miaoya ofrece una experiencia intensa, pero no necesariamente un reemplazo permanente; si un Agent trabaja de forma estable, es más probable que se convierta en una puerta de un solo sentido. \\
Lugar en la mente y retención & ¿Con qué palabra lo recordará el usuario? ¿Se usa con suficiente frecuencia? & PPT es demasiado estrecho; creación y expresión es más amplio; un amigo de IA, en cambio, requiere una sensación de relación a largo plazo. \\
\bottomrule
\end{longtable}

\begin{importantbox}{Del éxito viral al producto de largo plazo}
Un éxito viral suele demostrar que la tecnología, la difusión y la necesidad coincidieron una vez; un producto de largo plazo debe demostrar además lugar en la mente del usuario, frecuencia, retención, recompra y evolución de capacidades. Los productos de IA deben cuidarse en especial de «haber demostrado solo el wow, sin demostrar la puerta de un solo sentido».
\end{importantbox}

\subsection{Por qué «yo mismo lo usaría» sigue siendo importante}

Zhang Yueguang dice que le cuesta juzgar un producto que él mismo no usaría en absoluto. Este criterio no es aplicable a todos los gerentes de producto, pero es especialmente importante en la etapa temprana de la IA, porque muchos indicadores todavía no se estabilizan y los usuarios tampoco saben explicar con claridad lo que quieren. Que el propio fundador realmente no pueda prescindir de esta nueva solución suele ser la primera señal cualitativa.

\begin{warningbox}{El uso propio no es una validación universal}
El uso propio permite captar intuiciones tempranas, pero no sustituye la segmentación de usuarios, los escenarios reales ni la validación comercial. Sirve como señal de arranque, no como prueba definitiva.
\end{warningbox}

\subsection{Resumen de la sección}

La evaluación de productos de IA debe considerar a la vez la tecnología, la función de efecto, la puerta de un solo sentido y el lugar en la mente del usuario. Miaoya demostró que Zhang Yueguang puede crear un éxito viral; Dokie debe demostrar que puede construir un Agent correcto a largo plazo.
\section{Organización de productos de IA: Taste y trabajo en dos etapas}

Cuando el producto pasa del diseño de flujos al diseño de contexto, la forma de organizar el trabajo también tiene que cambiar. Zhang Yueguang (张月光) dice que los límites de las decisiones se vuelven difusos, que Taste adquiere importancia y que el trabajo deja de ser lineal para organizarse en dos etapas. Aquí Taste no es un gusto estético subjetivo, sino la capacidad de juzgar si la salida del modelo es buena o mala, dónde están los límites de la experiencia de usuario y si una dirección de producto merece la apuesta.

\begin{figure}[H]
\centering
\includegraphics[width=0.96\textwidth]{ai-product-org-two-stage.png}
\caption{Organización de productos de IA: el trabajo lineal se convierte en trabajo en dos etapas.}
\end{figure}

\begin{knowledgebox}{Lectura de la figura: primero validar la capacidad, después empaquetar el producto}
La primera etapa de la figura es la viabilidad del modelo: context, prompt, herramientas, combinación de modelos, evaluación y juicio de taste; solo la segunda etapa es el empaquetado del producto: flujos, UI, modelo de negocio y lugar en la mente del usuario. El flujo de trabajo lineal tradicional coloca la definición de requisitos al principio, mientras que un equipo AI Native primero debe confirmar que la capacidad del modelo existe realmente.
\end{knowledgebox}

\subsection{Por qué importa Taste}

La salida de la IA es incierta, y la entrada del usuario también lo es. El gerente de producto no puede limitarse a decir «aquí va un botón»; tiene que juzgar si la salida es lo bastante buena, si es controlable, si hay usuarios dispuestos a seguir usándola y si puede convertirse en una puerta de un solo sentido. Este juicio a menudo no puede sustituirse por completo con métricas, porque un producto en su etapa temprana todavía no tiene métricas estables.

\begin{knowledgebox}{Términos clave: organización de productos de IA}
\begin{tabularx}{\textwidth}{p{0.24\textwidth}p{0.32\textwidth}X}
\toprule
Término & Explicación breve & Papel en este episodio \\
\midrule
Taste & Capacidad de juzgar la calidad de la salida, la dirección del producto y la experiencia de usuario & Capacidad central cuando los límites de las decisiones son difusos. \\
Context Design & Diseñar la información, las herramientas y las restricciones visibles para el modelo & Capa intermedia de los productos AI Native. \\
Prompt Practice & Método para construir la entrada y la descripción de la tarea & Es solo una parte de context design. \\
Two-stage Work & Primero validar la capacidad del modelo, después empaquetar el flujo del producto & Sustituye la colaboración lineal tradicional. \\
Evaluation & Juzgar si la salida es estable, utilizable e iterable & Conecta taste con la ingeniería. \\
\bottomrule
\end{tabularx}
\end{knowledgebox}

\subsection{Resumen de la sección}

Lo esencial en la organización de productos de IA es que producto, diseño, ingeniería y capacidad del modelo trabajen dentro del mismo ciclo de retroalimentación. Sin taste, el equipo se deja deslumbrar fácilmente por una demo; sin validación de ingeniería, taste se vuelve pura especulación.
\section{Segundo emprendimiento: población de IA, Will y Skill}

Una vez expuestos los criterios para seleccionar productos, la discusión pasa a un juicio más general: qué está creando realmente esta generación de IA. La respuesta de Zhang Yueguang (张月光) no es «más servicios», sino «población de IA».

Cuando emprendió por segunda vez, Zhang Yueguang partía de un juicio más ambicioso: esta generación de tecnología de IA no crea servicios de IA, sino población de IA. Por población de IA se entiende la entrada en las redes humanas de una gran cantidad de individuos de IA capaces de interactuar, ejecutar y colaborar. Este juicio explica por qué trabaja al mismo tiempo en un juego otome con IA y en un Agent: el primero explora individuos de IA en el ámbito emocional y de contenido; el segundo, individuos de IA en el ámbito de la productividad.

\begin{figure}[H]
\centering
\includegraphics[width=0.94\textwidth]{ai-population-will-skill.png}
\caption{Población de IA, Will y Skill: la persona se encarga de los objetivos y la IA, de la ejecución.}
\end{figure}

\begin{knowledgebox}{Lectura de la figura: diferencia entre servicio de IA y población de IA}
Si la IA es un servicio, el usuario invoca funciones; si la IA es población, el usuario colabora con una entidad capaz de hacer cosas. En la figura, Will designa los objetivos, la intención, la elección y el gusto; Skill designa la ejecución, la generación, la invocación de herramientas y la superación de los límites de la capacidad individual.
\end{knowledgebox}

\subsection{La persona se encarga de la Will y la IA, de la Skill}

La sección anterior presentó la IA como «población»; esta sección desglosa la división del trabajo entre la persona y la IA. Will / Skill es un marco de producto breve, pero muy práctico.

En la entrevista, afirma que la persona debe encargarse de la Will y la IA, de la Skill. Esta frase sirve muy bien para explicar los productos de tipo Agent: el usuario no necesita dominar todas las habilidades, pero sí debe saber qué quiere, qué resultado es bueno, cuándo continuar y cuándo detenerse. La IA, por su parte, asume la ejecución, la invocación de herramientas y la generación.

\begin{importantbox}{División del trabajo entre Will y Skill}
Cuanto más potente es un producto de IA, menos debe excluir por completo a la persona. El valor de la persona se desplaza hacia los objetivos, el juicio, el gusto y la elección; el valor de la IA se desplaza hacia la ejecución, la generación y la ampliación de capacidades.
\end{importantbox}

\subsection{Resumen de la sección}

La población de IA no es un eslogan, sino un cambio en el posicionamiento del producto: el usuario no solo usa funciones, sino que colabora con individuos de IA. La división del trabajo entre Will y Skill es el marco central con el que Zhang Yueguang entiende los Agents y los amigos de IA.
\section{La perspectiva del CEO: el periodo de incertidumbre, la energía mental y el All In}

Otro aspecto importante de este episodio es cómo un emprendedor gestiona el periodo de exploración. Muyan Zhiyu (沐言智语) obtuvo financiamiento sin contratiempos, pero Zhang Yueguang (张月光) repitió varias veces que el dinero no es lo más importante: lo es más la energía mental del equipo y la suya propia. Una exploración sin resultados desgasta al equipo; si el fundador cambia de dirección una y otra vez sin gestionar las expectativas, el equipo termina por derrumbarse.

\subsection{Cómo comunicarse en el periodo de incertidumbre}

Describe los principios que sigue en el periodo de incertidumbre: atenerse a los hechos y no presentar como dirección segura una idea en la que no tiene confianza; decirle al equipo que todavía no se ha hecho all in; permitir que otros propongan mejores ideas; y probar con errores de bajo costo para cuidar la energía mental del equipo. Esta forma de actuar coincide con su método de exploración de producto: cuando hay incertidumbre, se controlan las variables; solo cuando hay certeza se concentran los recursos.

\begin{knowledgebox}{Gestión organizacional en el periodo de exploración}
En un emprendimiento de aplicaciones de IA es fácil caer en un estado en el que «cada dirección parece una oportunidad». En el periodo de exploración hay que proteger el efectivo y, más aún, la energía mental del equipo; antes de hacer all in, hay que decirle claramente al equipo que se trata solo de una prueba, no de la batalla definitiva.
\end{knowledgebox}

\subsection{Por qué con Dokie sí se puede hacer All In}

Zhang Yueguang cuenta que, cuando pensó en lo que hay detrás de Dokie ---crear un amigo de IA que rebase los límites de sus propias capacidades---, aceptaba incluso perder el dinero de la empresa en ese proyecto. Es un criterio exigente: no se trata de si un producto aislado de PPT puede volverse un éxito, sino de si la gran tesis que hay detrás merece que la empresa apueste por ella.

\begin{importantbox}{La decisión de hacer All In}
El all in no debe surgir de la ansiedad, de la presión del financiamiento ni de perseguir la moda, sino de la convicción del fundador sobre una gran tesis: aun si fracasa, sabe por qué valía la pena perder en ella.
\end{importantbox}

\subsection{Resumen de la sección}

Una empresa de aplicaciones de IA no solo tiene que encontrar un producto, también tiene que gestionar el periodo de exploración. El all in de Zhang Yueguang con Dokie es, en realidad, una apuesta por «un amigo de IA que rebase los límites de sus capacidades», y no por la herramienta puntual de PPT.
\section{Xingmian (星眠) y Dokie: dos líneas de producto}

Ya se explicó la decisión emprendedora de apostarlo todo (all in); ahora hay que ver en qué productos concretos se materializa esa decisión. Xingmian y Dokie representan dos rutas: el contenido de compañía y el Agent de productividad.

Muyan Zhiyu (沐言智语) tiene hoy dos líneas principales: el juego otome con IA Xingmian y el producto Agent Dokie. El primero continúa el interés de largo plazo de Zhang Yueguang (张月光) por Live2D, la compañía y el contenido dirigido al público femenino; el segundo es la dirección de productividad en la que decidió apostarlo todo (all in) tras un año de exploración.

\begin{figure}[H]
\centering
\includegraphics[width=0.94\textwidth]{xingmian-dokie-two-bets.png}
\caption{Las dos líneas de producto de Muyan Zhiyu: Xingmian y Dokie.}
\end{figure}

\begin{knowledgebox}{Lectura de la figura: uno hace contenido, el otro amplía los límites de capacidad}
Xingmian pone a prueba el contenido y la compañía con IA; lo que importa ahí es el contenido, los personajes, la frecuencia y la interacción. Dokie entra por las PPT y la expresión creativa, y su objetivo es que el usuario supere los límites de su capacidad. Ninguno de los dos es una simple demostración técnica: ambos buscan una puerta de entrada con la que el usuario interactúe, a la que recuerde y por la que pague a largo plazo.
\end{knowledgebox}

\subsection{Por qué los juegos con IA son un espacio estrecho}

Zhang Yueguang considera que el espacio viable de los juegos con IA es más estrecho de lo que mucha gente cree. La generación de video en tiempo real, los juegos cinematográficos de interacción infinita y las tramas dinámicas en mundo abierto todavía son difíciles de llevar a la práctica de forma estable con la tecnología actual. Donde la IA aporta valor con más facilidad es en la conversación de compañía y en la relación con los personajes; por eso confía más en los juegos otome con IA que en los juegos con IA en general.

\begin{warningbox}{Un juego con IA no es cualquier juego más IA}
En esencia, los juegos siguen siendo una industria de contenido. La IA puede cambiar ciertas experiencias de contenido, pero no resuelve por sí sola la trama, los costos, la estabilidad ni la retención de usuarios. Los escenarios viables suelen ser más estrechos que las grandes visiones.
\end{warningbox}

\subsection{Por qué Dokie entra por las PPT}

Dokie entra por las PPT, pero Zhang Yueguang insiste en que no será para siempre una herramienta de PPT. Las PPT son una puerta de entrada a la creación y la expresión; el objetivo real es «un amigo de IA que me ayude a superar los límites de mi capacidad». Esto explica por qué está dispuesto a apostarlo todo (all in) por el Agent y no a poner todo el capital en aquellas exploraciones pequeñas de antes.

\subsection{Resumen de la sección}

Xingmian y Dokie representan dos formas de la población de IA: el individuo de IA en la compañía y el contenido, y el individuo de IA en la productividad y la creación. El primero tiene que controlar el costo del contenido y los límites de la compañía; el segundo tiene que demostrar que de verdad se superan los límites de capacidad.
\section{Agent: retroalimentación de ciclo corto y alta frecuencia, y el LM como núcleo}

Después de hablar de Xingmian (星眠) y Dokie, la atención pasa naturalmente al criterio sobre Agent que hay detrás de Dokie. Esta sección divide el know-how de Zhang Yueguang (张月光) sobre Agent en dos partes: la ruta de producto y el núcleo técnico.

La segunda mitad de la entrevista entra en el tema de Agent. Zhang Yueguang ofrece dos know-how. Primero, no hay que ver al Agent solo como una herramienta para tareas largas y fuera de línea: cuanto más larga es la tarea, más inestable se vuelve; quizá convenga construir ciclos más cortos, más frecuentes y más cercanos a la retroalimentación del usuario. Segundo, al menos hoy, el LM sigue siendo el núcleo intelectual del Agent, porque es capaz de entender la intención del usuario, descomponer tareas, invocar herramientas y llamar a otros modelos.

\begin{figure}[H]
\centering
\includegraphics[width=0.96\textwidth]{agent-short-loop.png}
\caption{Ruta de retroalimentación de ciclo corto y alta frecuencia para productos de Agent.}
\end{figure}

\begin{knowledgebox}{Lectura de la figura: el ciclo corto y de alta frecuencia no significa poca capacidad, sino un producto controlable}
En la figura, la intención del usuario entra al núcleo LM; el LM invoca herramientas y modelos para generar productos, y el usuario vuelve a dar retroalimentación con alta frecuencia. En comparación con meter al Agent en una tarea larga y esperar el resultado final, el ciclo corto y de alta frecuencia facilita construir confianza, reduce la ansiedad por la latencia e integra el juicio del usuario en el ciclo cerrado.
\end{knowledgebox}

\subsection{Por qué el LM sigue siendo el núcleo}

Zhang Yueguang no está muy de acuerdo con la afirmación de que «el lenguaje no puede llegar al punto final de la inteligencia»; incluso considera que el lenguaje es la esencia de la inteligencia. Esta postura forma un contraste interesante con la visión de los modelos del mundo de Xie Saining (谢赛宁) en EP133. Zhang Yueguang habla desde la perspectiva del producto y del Agent: la intención del usuario se expresa principalmente con lenguaje, la invocación de herramientas requiere planificación en lenguaje y los productos del trabajo son, en gran medida, lenguaje o contenido estructurado; por eso el LM es hoy el núcleo rector más natural.

\begin{warningbox}{Aquí hay una divergencia valiosa}
Xie Saining enfatiza que el lenguaje no puede sustituir al mundo físico; Zhang Yueguang enfatiza que el lenguaje es el núcleo de inteligencia del Agent. Ambas posturas no se contraponen necesariamente: la inteligencia del mundo físico necesita ir más allá del lenguaje, mientras que los Agent de productividad dependen, en esta etapa, en gran medida del lenguaje para organizar las tareas.
\end{warningbox}

\subsection{Resumen de la sección}

Los productos de Agent no compiten solo en benchmark de tareas largas. Los ciclos cortos, la alta frecuencia, la baja latencia, la retroalimentación del usuario y los resultados editables pueden formar antes una experiencia de producto real. El LM sigue siendo, por ahora, el cerebro central del Agent.
\section{Plataformas, herramientas y la mente del usuario}

Una vez que el Agent puede hacer cosas, falta responder qué es en la mente del usuario. Esta pregunta define los límites del producto: si es una herramienta, una plataforma o un nuevo punto de entrada.

Zhang Yueguang (张月光) también habló de las oportunidades para nuevas plataformas. Según él, una nueva plataforma de contenido necesita un nuevo medio y una nueva forma de interacción. Sora no necesariamente se convertirá en plataforma, porque podría seguir siendo un formato de contenido de video corto. ChatGPT, en cambio, sí podría convertirse en plataforma, porque creó un nuevo punto de entrada para la interacción. En el caso de los Agent, términos como vertical/general o modelo/aplicación son demasiado gruesos: lo que realmente importa es con qué idea te recuerda el usuario.

\begin{figure}[H]
\centering
\includegraphics[width=0.94\textwidth]{platform-tool-agent-positioning.png}
\caption{Oportunidades para nuevas plataformas: productos de tipo herramienta, productos de tipo plataforma y el Agent como punto de entrada.}
\end{figure}

\begin{knowledgebox}{Lectura de la figura: la mente del usuario define los límites}
Una herramienta resuelve tareas concretas; una plataforma crea relaciones y un ecosistema. Un Agent como punto de entrada debe lograr que el usuario asocie la marca con una necesidad que tiene con frecuencia. Si el alcance es demasiado estrecho, se usa poco; si es demasiado amplio, «poder hacer de todo equivale a no poder hacer nada». Dokie se posiciona como un Agent de creación y expresión, no como una mezcla de todas las tareas posibles.
\end{knowledgebox}

\subsection{Los límites de la creación y la expresión}

Zhang Yueguang considera que las presentaciones (PPT), los documentos, la expresión y la creación pueden reunirse en una idea de nivel superior en la mente del usuario: ayudarlo a crear y a expresarse mejor. En cambio, tareas como recopilar información, llenar formularios o procesar datos no necesariamente se integrarán con un Agent de creación. Es un criterio de producto: los límites no se trazan según las capacidades del modelo, sino según la mente del usuario y los escenarios de uso.

\begin{importantbox}{El problema del posicionamiento de un Agent}
Un Agent no es mejor por ser más general. El producto debe encontrar en la mente del usuario un límite que este quiera recordar, que se use con frecuencia y que sea lo bastante amplio sin volverse vago.
\end{importantbox}

\subsection{Resumen de la sección}

La competencia entre productos de Agent no es solo una competencia de capacidades de los modelos, sino también por un lugar en la mente del usuario. Que pase de ser una herramienta a ser un punto de entrada depende de que tenga un límite de uso frecuente, una relación con el usuario y capacidad de seguir ampliándose.
\section{Conexión con los episodios anteriores: Agent, Context y lenguaje}

Este episodio forma, junto con los anteriores, una línea principal muy clara en la capa de aplicación. EP139 trata la historia técnica de los Agent y enfatiza la llamada a herramientas, la disolución de fronteras y la irradiación social; EP135 trata Elys y los dobles digitales, y enfatiza la obtención y el flujo de context y la red personificada; EP133, desde la perspectiva de los modelos del mundo, recuerda que más allá del lenguaje existe el mundo físico. El lugar de EP130 consiste en devolver estos conceptos técnicos a la mesa de trabajo del gerente de producto.

\subsection{Una misma palabra, significados en distintos niveles}

\begin{longtable}{p{0.22\textwidth}p{0.34\textwidth}p{0.35\textwidth}}
\toprule
Concepto & Énfasis en los episodios anteriores & Significado de producto en EP130 \\
\midrule
Agent & Llamada a herramientas, ejecución de tareas, cambios en los límites del sistema & Entra por la expresión en PPT y la creación, y se convierte en un objeto de productividad interactivo. \\
Context & Información personal, redes sociales, flujo de datos de los dobles digitales & El gerente de producto debe diseñar el contexto visible para el modelo, no solo dibujar flujos. \\
Lenguaje & En EP133 se discute como insuficiente para abarcar el mundo físico & En EP130 se considera el núcleo de gobierno del Agent y la puerta de entrada de la intención del usuario. \\
Amigo de IA & EP135 se inclina más hacia las relaciones y las redes sociales & EP130 enfatiza ofrecer primero valor práctico y, poco a poco, generar la sensación de amistad. \\
\bottomrule
\end{longtable}

\begin{knowledgebox}{Por qué importa esta conexión}
Si se mira solo EP130, parece una metodología de producto; puesto dentro de la serie completa, responde a «cómo un Agent se convierte en un producto que la gente común usa a largo plazo». La historia técnica, el flujo de context, los límites del lenguaje y la percepción del producto son distintos niveles de un mismo problema.
\end{knowledgebox}

\subsection{Resumen de la sección}

EP130 lleva al Agent del relato técnico de vuelta al relato de producto: no solo pregunta si el modelo puede hacerlo, sino también por qué el usuario lo recuerda, por qué ya no puede volver atrás y por qué está dispuesto a dejarlo entrar en su trabajo y en su red de relaciones.
\section{Términos clave: índice de palabras clave de este episodio}

\begin{longtable}{p{0.24\textwidth}p{0.32\textwidth}p{0.35\textwidth}}
\toprule
Término & Explicación en una frase & Papel en este episodio \\
\midrule
AI Native & El paradigma del producto lo reescribe la IA; no se trata solo de usar IA & Concepto central de todo el episodio. \\
Diseño de flujos & Rutas, botones, entradas y salidas predefinidos & Paradigma por defecto de los productos de internet. \\
Diseño del contexto & Diseñar la información, las herramientas y las restricciones visibles para el modelo & Trabajo central de un producto AI Native. \\
One Way Door & Una vez que el usuario adopta la nueva solución, no quiere volver a la anterior & Criterio para identificar un producto acertado a largo plazo. \\
Taste & Capacidad de juzgar la salida del modelo, la dirección del producto y la experiencia del usuario & Capacidad central de una organización de productos de IA. \\
Población de IA & Individuos de IA capaces de interactuar, ejecutar y colaborar se integran a la red humana & Visión macro de Zhang Yueguang (张月光) sobre esta generación de IA. \\
Will / Skill & La persona se encarga de la intención y la elección; la IA, de la capacidad de ejecución & División del trabajo entre personas y máquinas en los productos Agent. \\
Xingmian (星眠) & Juego otome con IA de Muyan Zhiyu (沐言智语) & Línea de contenido y compañía con IA. \\
Dokie & Producto Agent de Muyan Zhiyu & Línea de expresión creativa y límites de capacidad. \\
Núcleo LM & Usar un modelo de lenguaje para comprender, descomponer tareas e invocar herramientas & Núcleo intelectual de los Agent actuales. \\
\bottomrule
\end{longtable}

\subsection{Resumen de la sección}

En conjunto, estos términos forman una metodología de productos de IA: del flujo al contexto, del servicio a la población de IA, de la herramienta a la puerta de entrada del Agent, del wow moment a la One Way Door.

\section{Síntesis y ampliación}

\subsection{Conclusiones centrales}

\begin{enumerate}
\item Miaoya (妙鸭) es un producto de internet que aprovecha muy bien las capacidades de la IA, pero no es AI Native en el sentido que Zhang Yueguang (张月光) definió después.
\item El cambio decisivo de AI Native es pasar del diseño de flujos al diseño del contexto: el gerente de producto debe organizar la información, las herramientas y la retroalimentación que el modelo puede ver.
\item One Way Door es una señal importante para reconocer un producto acertado a largo plazo: una vez que el usuario migra, ya no quiere volver a la solución anterior.
\item La organización de un producto de IA deja de ser un trabajo lineal y pasa a tener dos etapas: primero se validan las capacidades del modelo y el context; después se encapsulan el flujo y la experiencia.
\item Zhang Yueguang entiende esta generación de IA como la creación de una «población de IA», y no solo como la creación de servicios de IA.
\item La división Will / Skill indica que la persona se encarga de los objetivos, el juicio y el gusto, y la IA, de la ejecución y de la ampliación de capacidades.
\item Xingmian (星眠) explora el contenido y la compañía con IA; Dokie explora un Agent de creación y expresión, y la superación de los límites de las capacidades del usuario.
\item Un producto de Agent puede empezar por una ruta de retroalimentación corta y frecuente, en lugar de perseguir solamente tareas fuera de línea de muy largo alcance.
\item El LM sigue siendo el núcleo de los Agent actuales, pero esa es una cuestión distinta de si la inteligencia en el mundo físico necesita ir más allá del lenguaje.
\end{enumerate}

\subsection{Preguntas abiertas}

Se dejan preguntas abiertas al final porque este episodio ofrece sobre todo metodología de producto y juicios de emprendimiento, no conclusiones que el mercado ya haya validado por completo. Las siguientes preguntas determinarán si Dokie y los productos de Agent similares pueden pasar de ser herramientas a ser puntos de entrada.

\begin{itemize}
\item ¿Los productos AI Native formarán una metodología estable, o cada tipo de producto tendrá que reinventar el context design?
\item ¿El mejor punto de entrada para un Agent es la creación y expresión, el procesamiento de información, el código, la automatización de oficina o un amigo de IA con personalidad?
\item ¿Pueden los juegos otome con IA convertirse en una categoría grande, o solo serán un camino estrecho dentro del contenido con IA?
\item ¿Cuánto tiempo podrá el LM seguir siendo el núcleo de los Agent? ¿Cambiarán este juicio los modelos multimodales?
\item ¿Los usuarios pagarán por un «amigo de IA», o solo por una «herramienta que trabaje de forma confiable»?
\end{itemize}

\subsection{Lecturas adicionales}

\begin{itemize}
\item EP139, historia técnica de los Agent: para entender la trayectoria de los Agent desde la invocación de herramientas hasta su alcance social.
\item EP135, Tristan / Elys: para entender los dobles digitales, el flujo del context y las redes sociales de IA.
\item EP113, entrevista a Yang Zhilin (杨植麟): Agentic LLM, el cerebro en una cubeta y «El comienzo del infinito».
\item Casos de coding agent como Cursor y Claude Code: para entender los productos One Way Door.
\item Materiales sobre juegos con IA, narrativa interactiva, Live2D y compañía virtual: para entender la línea de contenido de Xingmian.
\end{itemize}

\begin{importantbox}{El juicio final}
Lo que más vale llevarse de este episodio no es el juicio contraintuitivo de que «Miaoya no es AI Native», sino la metodología que lo sustenta: cuando tanto la entrada como la salida son abiertas, el producto deja de ser solo un flujo; se convierte en un sistema formado en conjunto por el contexto, el modelo, las herramientas, la retroalimentación, el gusto y la percepción del usuario.
\end{importantbox}

\end{document}
